\documentclass[11pt, a4paper, oneside]{article}
\usepackage[ngerman]{babel}
\usepackage{worksheet}

\title{Webentwicklung: CSS}
\author{L. Bung}
\subject{SAE}
\class{E3FI}

\begin{document}
	\maketitle
	
	\begin{hint}{Was ist CSS?}
		CSS (Cascading Style Sheets) ist eine Deklarationssprache für Stildefinitionen von HTML-Dokumenten.
		Während HTML die Struktur bereitstellt, übernimmt CSS die visuelle Gestaltung wie Farben, Abstände, Layouts und Schriftarten.
		
		Jeder Block besteht dabei aus einem \emph{Selektor} mit geschweiften Klammern.
		Dazwischen stehen mit Semikolon beendete Schlüssel-Wert-Paare (\emph{Properties}), die bestimmte Designregeln festlegen.
		
		\begin{lstlisting}[language=CSS, numbers=none]
h1 { //Selektor: Alle h1-Überschriften werden ausgewählt
	color: red; //Setze die Textfarbe auf rot.
}
		\end{lstlisting}
	\end{hint}
	
	
	\partnertask{Adressierung von Elementen über CSS-Selektoren}
	
	CSS-Regeln werden mit Hilfe von Selektoren an bestimmte HTML-Elemente gebunden.
	Sehen Sie sich den folgenden HTML-Ausschnitt an:
	
	\begin{lstlisting}[language=HTML]
<h1 class="title">Willkommen</h1>
<p class="text-highlight">Dieser Text ist hervorgehoben.</p>
<p id="special-note">Eine spezielle Anmerkung.</p>
<div>Ein gewöhnlicher Bereich.</div>
	\end{lstlisting}
	
	\begin{enumerate}[label=\alph*)]
		\item Schreiben Sie die passenden CSS-Selektoren auf, um:
		\begin{itemize}
			\item Alle \texttt{<p>}-Elemente auszuwählen.
			\item Das Element mit der Klasse \texttt{text-highlight} auszuwählen.
			\item Das Element mit der eindeutigen ID \texttt{special-note} auszuwählen.
		\end{itemize}
		\item Vergleichen Sie Ihre Erläuterung mit Ihrem Sitznachbarn und erklären Sie den fachlichen Unterschied zwischen einer Klasse (\texttt{class}) und einer ID (\texttt{id}).
	\end{enumerate}
	
	\lines[4cm]
	
	
	\task{Properties}
	
	CSS-Anweisungen bestehen aus Eigenschaft-Wert-Paaren (\texttt{property: value;}).
	Erarbeiten Sie sich anhand der Dokumentation (z.\,B. MDN Web Docs) grundlegende Styling-Eigenschaften.
	
	\begin{enumerate}[label=\alph*)]
		\item Nennen Sie die gesuchte CSS-Eigenschaft:
		
		\begin{itemize}
			\itemsep10pt
			\item Schrifttyp/Schriftart festlegen: \hrulefill
			\item Textfarbe ändern: \hrulefill
			\item Innenabstand eines Elements angeben: \hrulefill
			\item Außenabstand eines Elements angeben: \hrulefill
			\item Rahmen um ein Element zeichnen: \hrulefill
		\end{itemize}
		
		\item Recherchieren Sie selbstständig den Unterschied zwischen \texttt{padding} und \texttt{margin}.
		Zeichnen Sie eine kleine Skizze zur Veranschaulichung.
	\end{enumerate}
	
	\boxarea[13cm]
	
	
	\singletask{Möglichkeiten der CSS-Einbindung}
	
	Um CSS-Regeln auf eine Website anzuwenden, gibt es drei grundlegende Verfahren.
	Betrachten Sie die verschiedenen Möglichkeiten der Einbindung.
	
	\begin{enumerate}[label=\alph*)]
		\item Beschreiben Sie die folgenden drei Ansätze zur Einbindung von CSS kurz in eigenen Worten.
		\begin{itemize}
			\item Inline-Style (über das HTML-\texttt{style}-Attribut)
			\item Interner Style-Block (über das \texttt{<style>}-Element im \texttt{<head>})
			\item Externe CSS-Datei (über das \texttt{<link>}-Element)
		\end{itemize}
		\item Erläutern Sie, warum in der professionellen Softwareentwicklung fast ausschließlich externe CSS-Dateien verwendet werden.
	\end{enumerate}
	
	\lines[4cm]
	
	
	\partnertask{Praktische Gestaltung der HTML-Seite}
	
	Greifen Sie auf Ihren HTML-Steckbrief aus der letzten Stunde zurück.
	Erstellen Sie eine externe CSS-Datei namens \texttt{style.css} und binden Sie diese in Ihre HTML-Datei ein.
	
	\begin{enumerate}[label=\alph*)]
		\item Gestalten Sie die Seite nach folgenden Vorgaben:
		\begin{itemize}
			\item Die Hauptüberschrift (\texttt{<h1>}) soll eine dunkelblaue Schriftfarbe und eine serifenlose Schriftart erhalten.
			\item Der Hintergrund des gesamten Dokuments (\texttt{<body>}) soll in einem hellen Grauton erscheinen.
			\item Das Firmenlogo (\texttt{<img>}) soll eine maximale Breite von $200\,\text{px}$ und einen abgerundeten Rahmen besitzen.
			\item Hervorgehobene Absätze sollen durch eine eigene Klasse mit einem linken bunten Rand markiert werden.
		\end{itemize}
		\item Sprechen Sie sich mit Ihrem Nebensitzer ab, sodass Sie beide dieselben Klassen verwenden.
		Tauschen Sie anschließend Ihre erstellten CSS-Dateien aus und testen Sie, ob alles wie geplant aussieht.
	\end{enumerate}
	
	
	\partnertask{Pseudoklassen und Media Queries}
	
	Webseiten müssen sowohl auf Mausinteraktionen reagieren als auch auf unterschiedlichen Bildschirmgrößen (vom Smartphone bis zum Desktop-Monitor) optimal dargestellt werden.
	Dafür stellt CSS sogenannte Pseudoklassen und Media Queries (\texttt{@media}) bereit.
	
	Sehen Sie sich den folgenden CSS-Ausschnitt an:
	
	\begin{lstlisting}[language=CSS]
a {
	color: blue;
	text-decoration: none;
}
a:hover {
	color: orange;
	text-decoration: underline;
}
@media (max-width: 768px) {
	.navigation {
		display: flex;
		flex-direction: column;
	}
}
	\end{lstlisting}
	
	\begin{enumerate}[label=\alph*)]
		\item Erklären Sie die Funktionsweise der Pseudoklasse \texttt{:hover} im obigen Beispiel.
		Nennen Sie zwei weitere gebräuchliche Pseudoklassen für Links oder Formularelemente.
		\item Analysieren Sie den \texttt{@media}-Block.
		Beschreiben Sie kurz, unter welcher konkreten Bedingung die darin enthaltene CSS-Regel angewendet wird und welche Auswirkung dies auf das Layout der Klasse \texttt{.navigation} hat.
		\item Diskutieren Sie kurz mit Ihrem Partner, warum das Prinzip des \textit{Responsive Web Design} in der modernen Softwareentwicklung unverzichtbar ist.
		\item Überarbeiten Sie die in Aufgabe 4 erstellte CSS-Datei.
		Nutzen Sie Media Queries, um Designregeln für mobile Endgeräte festzulegen\footnotemark.
	\end{enumerate}
	
	\footnotetext{\emph{Hinweis}: In der Entwicklerkonsole des Browsers gibt es die Option, ein Mobilgerät zu simulieren.}
	
	
	\pagebreak
	
	\bonustask{Das Kaskaden-Prinzip und Spezifität}
	
	Was passiert, wenn mehrere CSS-Regeln auf dasselbe HTML-Element zutreffen?
	CSS nutzt Regeln der \textit{Spezifität} (Specificity), um Zielkonflikte zu lösen.
	
	Gegeben ist folgender Code:
	
	\begin{lstlisting}[language=HTML]
<p id="first-para" class="highlight" style="color: green;">Welche Farbe habe ich?</p>
	\end{lstlisting}

	\begin{lstlisting}[language=CSS]
p { color: red; }
#first-para { color: blue; }
.highlight { color: orange; }
	\end{lstlisting}
			
	\begin{enumerate}[label=\alph*)]
		\item Bestimmen Sie, welche Farbe der Text letztendlich im Browser annimmt.
		\item Begründen Sie Ihre Entscheidung anhand der Rangfolge der Spezifität (Inline-Style vs. ID vs. Klasse vs. Tag-Selektor).
	\end{enumerate}
	
\end{document}